% Source: scan.pdf pages 1-2 and the top of page 3.
% Chapter 2 starts partway down page 3; that part is in ../ch02/notes.tex.
\documentclass[11pt]{article}
\usepackage[margin=1in]{geometry}
\usepackage{parskip}
\usepackage{amsmath,amssymb}
\usepackage{booktabs}
\usepackage{graphicx}
\usepackage{xcolor}
\usepackage{tikz}
\usepackage{pgfplots}
\pgfplotsset{compat=1.18}
\usepackage[hidelinks]{hyperref}

% Uncertain reading from the scan. Prints red until you fix it against the paper.
\newcommand{\unclear}[1]{\textcolor{red}{[#1?]}}

% Explanation added during processing, not written on the page.
\newenvironment{aside}{\begin{quote}\small\color{black!65}\textbf{Note.}\ }{\end{quote}}

\title{ISE 555 Reading Notes\\\large Chapter 1: Introduction to Operations Research}
\author{Luke Pepin}
\date{September 21, 2026}

\begin{document}
\maketitle

\section{Origins}

Operations research (OR) grew out of World War II. It has been described as
``the art of winning the war without actually fighting it'' (Arthur C.\ Clarke).
\begin{itemize}
    \item \textbf{British:} efficient ocean transport, effective bombing.
    \item \textbf{U.S.:} fighting patterns, planning, utilization of electronics.
\end{itemize}

\section{Definition}

\begin{quote}
``The application of methods of science to complex problems in the direction and
management of large systems of men, machines, materials and money in industry,
business and defence\ldots\ Develop a scientific model of the system, incorporating
measurements of factors such as chance and risk with which to predict and compare
the outcomes of alternative decisions, strategies or controls. Help management
determine its policy and actions scientifically.''
\end{quote}

\begin{aside}
This is the Operational Research Society (UK) definition, one of twelve definitions
the book lists in Section 1.3.
\end{aside}

\section{Applications}

\begin{itemize}
    \item Budgeting and investments
    \item Purchasing
    \item Production management
    \item Marketing
    \item Personnel management
    \item Research and development
\end{itemize}

\section{Phases of an OR Study}

\begin{enumerate}
    \item \textbf{Formulation of the problem:} objectives, variables, ranges.
    \item \textbf{Construction of a mathematical model:} iconic, analogue, or
          symbolic (mathematical).
    \item \textbf{Deriving a solution from the model:} optimize for the objective,
          plus sensitivity analysis.
    \item \textbf{Testing the model and solution:} the model's solution should
          perform best.
    \item \textbf{Implementing a solution:} be ready, since a solution can get
          outdated.
\end{enumerate}

\section{Scope}

A wide variety of industries apply OR: it applies anywhere there is a
\emph{strategy} to choose.

\section{Modelling in OR}

A model is ``an idealized representation of a real-life system.''

\begin{itemize}
    \item \textbf{Advantages:} concise, logical, systematic; enables high-powered math.
    \item \textbf{Disadvantages:} not to be considered absolute; verified only
          through experiment; \unclear{timely}.
\end{itemize}

\subsection*{Types of models}

\begin{itemize}
    \item \textbf{Iconic:} a pictorial representation of a system, scaled up or
          down for a purpose (a globe, an atom diagram, a map). Difficult to
          manipulate for experiments.
    \item \textbf{Analogue:} one set of properties is used to represent another.
          In a graph, distance represents a variable; in a pie graph, area
          represents a percentage or probability. Not very specific.
    \item \textbf{Mathematical (symbolic):} uses equations to predict behavior.
    \begin{enumerate}
        \item[(i)] \textbf{Deterministic:} parameters are assumed known and precise.
        % CHANGED: page had "(general)" after the non-deterministic definition.
        % Specific/general is a separate classification (below), not a synonym for non-deterministic.
        \item[(ii)] \textbf{Non-deterministic:} at least one parameter is random or fuzzy.
    \end{enumerate}
\end{itemize}

\subsection*{Specific vs.\ general models}

\begin{itemize}
    \item \textbf{Specific} models are either \textbf{static} or \textbf{dynamic},
          depending on whether a time factor is considered.
    \item \textbf{General} models are non-optimized.
\end{itemize}

\subsection*{General solution methods}

\begin{enumerate}
    \item[(i)] \textbf{Analytical methods:} basic math, differentials, calculus.
    \item[(ii)] \textbf{Numerical methods:} trial and error, machine learning,
          continuous improvement.
    \item[(iii)] \textbf{Monte Carlo technique:} random number table; plug in and find.
\end{enumerate}

\section{From Juan's Notes (September 14 and 16)}

\begin{itemize}
    \item Elementary calculus is not always enough: there can be several turning
          points, local vs.\ global minima and maxima, and endpoints to check.
    \item Standard problem formulation:
    \begin{enumerate}
        \item Decision and scope
        \item Given data
        \item Decision variables
        \item Objective
        \item Constraints
        \item Domain and assumptions
        \item Validate and interpret
    \end{enumerate}
\end{itemize}

\subsection*{Examples}

\textbf{4. Two-variable linear program}

\begin{minipage}[t]{0.45\linewidth}
\vspace{0pt}
Maximize: $f(x) = 2x_1 + x_2$

Subject to:
\begin{itemize}
    \item $x_1 + x_2 \le 1$
    \item $x_1, x_2 \ge 0$
\end{itemize}
\end{minipage}\hfill
\begin{minipage}[t]{0.45\linewidth}
\vspace{0pt}
\centering
\begin{tabular}{ccc}
    \toprule
    Point & Coordinates & Objective value \\
    \midrule
    $A$ & $(1, 0)$ & 2 \\
    $B$ & $(0, 1)$ & 1 \\
    $C$ & $(0, 0)$ & 0 \\
    \bottomrule
\end{tabular}
\end{minipage}

\medskip
\textbf{5. Cabinet manufacturing}

\textbf{6. Insurance workload allocation}

\begin{aside}
The numbers 4 to 6 match the section numbers in Juan's notes, where these
examples are worked in full: \texttt{Notes/Lectures/2026-09-14/J-notes.pdf}.
\end{aside}

\end{document}
